\section{Limitations and Future Research}

The empirical analysis is confined to Shariah-compliant equities in the BIST
Participation All Index (BIST XKTUM) over 2007--2022. The Turkish market has
a distinctive macroeconomic profile characterised by high inflation, significant
exchange rate volatility and a participation finance regulatory framework that
differs from those in Malaysia, the GCC countries or Indonesia. Consequently,
the relative ranking of models and the magnitude of portfolio performance metrics
may not transfer directly to other Islamic equity markets. Cross-market replication
using the Dow Jones Islamic Market Index, the FTSE Bursa Malaysia Hijrah Shariah
Index or GCC-listed Shariah-screened equities would allow researchers to assess
whether the superiority of XGBoost for prediction and Bi-LSTM for risk-adjusted
portfolio performance holds under different regulatory and macroeconomic conditions.

The study also relies on a single Shariah screening standard, namely the BIST XKTUM
classification issued under Turkish participation finance regulation. Different Shariah
supervisory boards apply varying financial ratio thresholds and business activity
screens, which means the composition of the investable universe --- and therefore the
covariance structure of the portfolio --- may differ when alternative screening
methodologies are applied. Future work could examine how screening methodology
affects the statistical properties of the investment universe and the subsequent
optimisation outcome, a question that is relevant for global Islamic fund managers
who operate across multiple jurisdictions.

\subsection{Portfolio optimisation framework}

The optimisation stage is restricted to LSTM, Bi-LSTM and GRU because these
architectures produce sequential return forecasts that fit naturally into the
rolling mean-variance framework. The paper therefore cannot directly compare
the portfolio-level value of XGBoost against the recurrent models, even though
XGBoost achieves the strongest prediction accuracy. Future research should explore
whether the point forecasts produced by gradient-boosting and ensemble models can
be converted into portfolio inputs --- for example, through quantile regression or
probabilistic forecasting --- and whether doing so produces portfolio outcomes
comparable to those of sequence-based models.

The optimisation constraints impose full investment and long-only weights, which
is appropriate for Shariah-compliant portfolios where short selling raises additional
jurisprudential concerns. However, the static nature of the framework means that
portfolio weights are not adjusted in response to structural breaks or volatility
regime changes within the sample period. A dynamic or regime-switching
optimisation approach, in which weights are revised at regular rebalancing
intervals conditional on updated model predictions, would better reflect the
adaptive management requirements of Islamic investment funds and is a natural
extension of the current framework.

\subsection{Benchmarking and comparative evaluation}

The portfolio evaluation compares LSTM, Bi-LSTM and GRU strategies against
each other using ending equity, annualised Sharpe ratio \citep{sharpe1966mutual}
and annualised volatility. The analysis does not include a comparison against
passive or heuristic benchmarks such as the equal-weight portfolio or the
BIST XKTUM index return. Without such benchmarks, it is not possible to
determine how much of the portfolio performance is attributable to the machine
learning forecasts and how much reflects the general upward trend in the index
over the sample period. Future research should include the equal-weight portfolio
\citep{taljaard2020equal, sen2023equal} and the Minimum Variance Portfolio as
reference points, since both are transparent strategies that are fully compatible
with Shariah-compliant constraints.

\subsection{Feature engineering and alternative inputs}

The feature set draws on price-derived variables including daily and 20-day returns,
rolling volatility and standard technical indicators. The models do not incorporate
macroeconomic variables such as exchange rates, commodity prices or inflation
proxies, nor do they use text-based sentiment signals from news or financial
disclosures. In Shariah-compliant equity markets, where firm eligibility depends
partly on balance-sheet ratios that change at reporting dates, fundamental accounting
data could also serve as an additional feature layer. Incorporating these data
sources would extend the information set available to the models and may
particularly benefit non-sequential models such as XGBoost, which already
demonstrates strong performance on structured tabular inputs.

\subsection{Model architecture}

The deep learning models in this study are implemented as recurrent architectures
without attention mechanisms. Transformer-based models, which replace recurrence
with self-attention and have produced strong results in financial forecasting tasks,
were not evaluated because their computational requirements at the scale of 163
stocks over 15 years would substantially increase GPU time beyond what was
feasible in this study. As GPU resources become more accessible, future work
should benchmark transformer architectures \citep{vaswani2017attention} against
the recurrent models evaluated here, particularly for longer lookback periods where
self-attention may capture dependencies that LSTM and GRU gates attenuate over time.

The hyperparameter search uses a predefined grid for each model class. More
data-efficient tuning strategies such as Bayesian optimisation
\citep{bergstra2012random} could identify strong configurations with fewer
evaluations, which would be valuable for researchers under tighter computational
constraints and would allow a wider architecture search in future Islamic
computational finance studies.

\subsection{Dimensionality reduction}

PCA is applied as a linear dimensionality reduction step before model training.
Because PCA assumes that the most informative directions in the data are those
with the highest variance, it may not preserve non-linear cluster structures that
exist among Shariah-compliant stocks with similar sector exposures or Shariah
ratio profiles. Non-linear alternatives such as UMAP or autoencoder-based
compression may better retain the complex relational structure of the 163-stock
universe, and their effect on downstream prediction accuracy and portfolio
construction is an open empirical question.
